% \documentclass[twocolumn]{article}
\documentclass[12pt]{article}
\usepackage{graphicx,amsmath,fancyhdr,fancyvrb,bm}
\usepackage[upint]{stix}
\title{Buborék dinamika és 2D akusztikus mez\H o szimuláció egyirányú csatolással}
\author{Gedei János Jonatán \\ Konzulens: Dr. Heged\H us Ferenc}
\date{2026 november}
\pagenumbering{arabic}     % Use Arabic numerals (1, 2, 3, ...)
\setcounter{page}{0}       % Start from page number 0
\thispagestyle{empty}      % Removes number on the current page (the first one)
\usepackage[a4paper,
            bindingoffset=0in,
            left=2.5cm,
            right=2.5cm,
            top=2.5cm,
            bottom=2.5cm,
            footskip=.25in]{geometry}
\pagestyle{fancy}
\fancyhf{} % clear all header and footer fields
% Put section title on the top (left or right depending on odd/even pages)
\fancyhead[L]{\nouppercase{\leftmark}}
% Page number at the bottom center
\fancyfoot[C]{\thepage}
% invert colors comment untill star to get a printable version
\usepackage{xcolor}
\definecolor{mybg}{RGB}{30,30,30}
\usepackage{background}
\backgroundsetup{
  scale=1,
  % color=black,
  % color=darkgray,
  % color=mybg,
  color=white,
  opacity=1,
  angle=0,
  position=current page.center, % center aligns full-page rectangle
  contents={\rule{\paperwidth}{\paperheight}}
}
% \color{white}
% star

% ----------------------------
% ---------- MACROS ----------
% ----------------------------
% numberless section title + add this to the header
\newcommand{\ms}[1]{%
  \newpage%
  \section*{Week #1, Lecture #1}%
  \markboth{Week #1, Lecture #1}{}%
  \setcounter{equation}{0}
  \setcounter{section}{#1}
  \setcounter{tocdepth}{1}
  \addcontentsline{toc}{section}{Week #1, Lecture #1}
}
% partial derivative, fraction
\newcommand{\pd}[2]{\frac{\partial #1}{\partial #2}}
% upper dagger
\newcommand{\adj}[0]{^\dag}
% braket, bra, ket notations
\newcommand{\braket}[2]{\left\langle #1 \middle| #2 \right\rangle}
\newcommand{\ketbra}[2]{\left| #1 \middle\rangle \middle\langle #2 \right|}
\newcommand{\bra}[1]{\left\langle #1 \right\vert}
\newcommand{\ket}[1]{\left\vert #1 \right\rangle}
% scalar product
\newcommand{\scp}[2]{\left\langle #1 , #2 \right\rangle}
% better sum
\newcommand{\summ}[2]{\sum_{#1}^{#2}}
% evaluated at
\newcommand{\at}[2]{\left. #1 \right\rvert_{#2}}
% array vertical column form stack exchange
\makeatletter
\renewcommand*\env@matrix[1][*\c@MaxMatrixCols c]{%
  \hskip -\arraycolsep
  \let\@ifnextchar\new@ifnextchar
  \array{#1}}
\makeatother
% set,where { #1 | #2 }
\newcommand{\setwhere}[2]{\begin{Bmatrix}[c|c] #1 & #2 \end{Bmatrix}}
% norm
\newcommand{\norm}[1]{\left\vert \left\vert #1 \right\vert \right\vert}
% reference equation
\newcommand{\refeq}[1]{equation (\ref{ #1 })}
% Hilbert-Schmidt scalar product
\newcommand{\scphs}[2]{\scp{#1}{#2}_{H-Sch}}
% quantum system state space (S^+_1(\mathcal{H}))
\newcommand{\qss}[0]{S^+_1(\mathcal{H})}
\newcommand{\qssi}[0]{S^+_1(\mathcal{H}_I)}
\newcommand{\qssii}[0]{S^+_1(\mathcal{H}_{II})}
\newcommand{\qssprod}[0]{S^+_1(\mathcal{H}_I\otimes\mathcal{H}_{II})}
\newcommand{\qssv}[0]{S(\mathcal{H})}
\newcommand{\qssiv}[0]{S(\mathcal{H}_I)}
\newcommand{\qssiiv}[0]{S(\mathcal{H}_{II})}
\newcommand{\qssprodv}[0]{S(\mathcal{H}_I\otimes\mathcal{H}_{II})}
% quantum system state space (S^+_1(  ))
\newcommand{\qssf}[1]{S^+_1(#1)}
% align* equation numbering
\newcommand{\nt}{\stepcounter{equation}\tag{\theequation}}
% in equation, often used words
\newcommand{\und}[0]{\text{ and }}
\newcommand{\oder}[0]{\text{ or }}
\newcommand{\tr}[1]{\text{ Tr}\left( #1 \right)}
\newcommand{\Span}[1]{\text{ Span}\left\{ #1 \right\}}
% \newcommand{\span}[1]{\text{ Span }\left( #1 \right) }
\newcommand{\carray}[1]{\begin{pmatrix} #1 \end{pmatrix}}
\newcommand{\lin}[1]{\text{Lin}\left( #1 \right)}
\newcommand{\subsectiontoc}[1]{%
  \subsection*{#1}%
  \addcontentsline{toc}{subsection}{#1}%
}
\newcommand{\subsubsectiontoc}[1]{%
  \subsubsection*{#1}%
  \addcontentsline{toc}{subsubsection}{#1}%
}
\newcommand{\paragraphtoc}[1]{%
  \paragraph{#1}%
  \addcontentsline{toc}{paragraph}{#1}%
}

\setlength{\parindent}{0pt} % removes paragraph indentation

\renewcommand{\contentsname}{Tartalom jegyzék}

% equation numbering as <#section.x>
\numberwithin{equation}{section}
% ------------------------------
% ---------- DOCUMENT ----------
% ------------------------------
\begin{document}

\maketitle
\begin{figure}[h]
    \centering
    \includegraphics[width=0.6\linewidth]{bme_logo.png}
\end{figure}
\begin{figure}[h]
    \centering
    \includegraphics[width=0.4\linewidth]{HDSlogo_szoveg.png}
\end{figure}

\thispagestyle{empty}

\clearpage
\thispagestyle{empty}
\null
\newpage

\section*{Absztrakt}
Ez a dolgozat  az akusztikus tér mozgó csomópontjainak hatását vizsgálja buborékok sugár és pozíció dinamikájára, numerikus módszerekkel.
Ezek az állapotváltozók jelent\H os szerepet játt-szanak a szonokémiai reaktorokban.
Korábbi munkánkat folytatva, két dimenziós szimulációt fejlesztettünk, amely az akusztikus hullámegyenletet oldja meg, GPU architektúrára illesztve, véges differenci módszerrel.
A szimulációkat el\H oször állóhullámok, majd a gerjesztés id\H ofügg\H o fáziseltolásával a nyomástér csomópont pozícióinak a mozgását reprodukáltuk.
Ezzel párhuza-mosan a rendszer CPU-ján egy irányú csatolással a Keller-Miksis és az impulzus egyenlet oldom meg, a buborék sugár és pozíció dinamikáját modellezve.
Végül pedig egyéb kutatások mérési eredményeivel összehasonlítottuk a szimulációk megoldásait.

\section*{Abstract}

This paper aims to examine the effects of moving nodes and anti-nodes of an acoustic field on bubble dynamics with changing radius and position with the use of numerical methods. These properties are essential for sonochemical reactors. Continuing our work, we developed a two dimensional acoustic simulation, solving the wave equation, mapped to a specific GPU architecture using finite difference method, for rapid field simulations. We started with simulating standing waves, then with an introduction of a time dependent phase shift for one of the transducer we reproduced the movement of the loction of the nodes and anti-nodes of the filed. Running in parallel on the systems CPU, with one way coupling to the Keller-Miksis and the momentum equations was added in order to model bubble radius and translation dynamics. Finaly our results were compared to researches measurement result, that examined bubble translation oscillations.

\section*{Köszönetnyilvánítás}

Szeretném hálámat kifejezni Nagy Dánielnek aki az ALPACA megoldó eredményeit szolgáltatta és Dr. Klapcsik Kálmánnak, aki a buborék dinamikai egyenleteivel segített.
Továbbá köszönettel tartozok konzulensemnek, Dr. Heged\H us Ferencnek, a dolgozat készítése során nyújtott szakmai támogatásáért.

This study was supported by the NVIDIA Corporation via the Academic Hardware Grants Program.


\newpage          

\setcounter{tocdepth}{4}
\tableofcontents
\newpage          

\section{Bevezetés}
\subsection{Szonokémia}

Miert is fontos a szonokemia?
Kemiai reaktorok iparban alkalmazott gepek, amik kornyezetet biztositanak kemiai reakciok vegbemenetelenek.
Egy gyakori pelda a Haber-Bosch folyamat, amely magas nyomason, es magas homersekleten 


\subsection{Miért kellett ezt a dolgozatot megírni}
\section{Modell}

\subsection{Fizikai modell => Matematikai modell}
Mi az a reaktor

Fizikai modell

\subsection{Numerikus módszerek}
fdm+rk45
haradware aware diszkretizáció

\subsection{Validáció}
ND's sol from ALPACE compared to my sol form gpu

\section{Akusztikus mez\H{o}}
\section{Eredmények}

\begin{figure}[h]
    \centering
    \includegraphics[width=0.5\linewidth]{placeholder.jpg}
    \caption{Oszilláció amplitudó, $|A_x|(\Delta f)$}
    \label{fig:1}
\end{figure}

\begin{figure}[h]
    \centering
    \includegraphics[width=0.5\linewidth]{placeholder.jpg}
	\caption{TF heatmap gain, $\at{G_x}{\Delta f}(R_0,P_a)$}
    \label{fig:1}
\end{figure}

\begin{figure}[h]
    \centering
    \includegraphics[width=0.5\linewidth]{placeholder.jpg}
	\caption{TF heatmap phase offset, $\at{\varphi_x}{\Delta f}(R_0,P_a)$}
    \label{fig:1}
\end{figure}


\section{compare to J's measurements?}
\section{Összefoglaló}

\section{Források}

\section{MI(AI) használat}

\end{document}
